\textbf{Contexto:} Nos últimos anos, o campo da inteligência artificial passou por um processo de expansão notável, tanto a nível acadêmico como industrial. Este crescimento pode ser observado de várias formas, incluindo o desenvolvimento de tecnologias mais complexas, o aumento do investimento, uma maior atenção dos meios de comunicação social e a expansão das suas áreas de aplicação. No entanto, este avanço deu origem a questões éticas que estão a tornar-se uma questão de crescente preocupação social, como o enviesamento de sistemas, aplicações danosas, entre outros.
\textbf{Objetivo:}  
Este trabalho teve como objetivo apresentar uma visão geral do estado atual das soluções práticas para o desenvolvimento ético de sistemas baseados em IA em todas as fases do ciclo de vida de um sistema de software. Pretendeu-se também desenvolver uma ferramenta que operacionalize a tradução de requisitos éticos de alto nível em histórias de usuário éticas.
\textbf{Método:} 
Para atingir este objetivo, este estudo baseou-se na metodologia \textit{Design Science Research}, que produziu três resultados até ao momento. Inicialmente, na fase de consciência do problema, foi atualizada a revisão sistemática da literatura desenvolvida por Cerqueira \cite{cerqueirahiccs}. Posteriormente, na fase de sugestão, a ferramenta proposta foi formulada como um sistema que traduz requisitos éticos em histórias de usuário éticas. Finalmente, na fase de desenvolvimento, foi implementada a primeira versão da ferramenta.
\textbf{Resultados:} No total foram identificados 38 estudos primários. Dentre estes estudos, a maior parte (63\%) propõe uma solução prática para facilitar a aplicação de ética em IA, mas ainda existe uma lacuna entre teoria e prática no que se diz sobre ética em IA. Além disso, foi também compilada uma lista com 26 dos principais princípios éticos que foram discutidos na literatura. Foi também proposta e desenvolvida a ferramenta \textit{US Translator}, que utiliza histórias de usuário, princípios éticos em IA e modelos de LLMs para promover a integração da ética em IA durante a fase de Engenharia de Requisitos.
\textbf{Conclusão:} Os resultados do processo de treinamento do modelo indicam que sua utilização é viável para possibilitar a prática de ética em IA durante os estágios iniciais do ciclo de vida do software. Consequentemente, os resultados preliminares indicam que a ferramenta facilita a integração da teoria e prática, preenchendo assim uma lacuna existente no atual conjunto de aplicações práticas.